\section{Conceptos básicos de distribución cuántica de claves de variable discreta}

Antes de comparar protocolos concretos, conviene fijar algunos conceptos básicos. Según \cite{gisin2002}, la distribución cuántica de claves, o QKD, no es un algoritmo de cifrado de mensajes, sino un mecanismo para que dos partes legítimas, normalmente llamadas Alicia y Bob, generen una clave secreta compartida. Como explican \cite{lo2014}, esa clave puede utilizarse después en un cifrado simétrico clásico; la función de QKD es, por tanto, resolver el problema de distribución de claves, no sustituir a todos los mecanismos de cifrado. La propuesta original de \cite{bennett1984} ya muestra que la seguridad del intercambio no se basa únicamente en la dificultad computacional de un problema matemático, sino en propiedades físicas de los sistemas cuánticos; en implementaciones reales, este razonamiento debe completarse con un modelo explícito de los dispositivos empleados \citep{scarani2009,lo2014}.

En los protocolos de variable discreta, la información se representa mediante cúbits. De acuerdo con \cite{nielsen2010}, un cúbit es un sistema cuántico de dos niveles que puede codificar los valores lógicos 0 y 1 mediante estados físicos concretos. En una implementación óptica, por ejemplo, esos estados pueden corresponder a polarizaciones, fases o intervalos temporales de pulsos luminosos \citep{gisin2002,scarani2009}. Por este motivo, a nivel abstracto resulta más riguroso hablar de cúbits, aunque en la práctica muchos sistemas los implementen mediante fotones o pulsos ópticos atenuados \citep{gisin2002,lo2014}.

Un protocolo DV-QKD combina normalmente dos canales. El canal cuántico transporta los cúbits o los sistemas físicos que los codifican, mientras que el canal clásico permite anunciar bases, descartar mediciones incompatibles, estimar errores, corregir discrepancias y aplicar amplificación de privacidad. Este canal clásico no necesita ser secreto, pero sí autenticado; de lo contrario, un atacante podría realizar un ataque de intermediario y hacerse pasar por una de las partes \citep{bennett1984, gisin2002,scarani2009,lo2014}.

Otro concepto central es el de base de preparación y medida. Como explican \cite{nielsen2010}, una base puede entenderse como un conjunto de estados distinguibles entre sí dentro de una misma medición. En la propuesta de \cite{bennett1984}, Alicia prepara cada cúbit en una de dos bases incompatibles, por ejemplo una base rectilínea o una base diagonal, y Bob elige la base con la que realiza la medición. Cuando ambos usan la misma base, el bit puede recuperarse con alta probabilidad en ausencia de ruido; cuando Bob mide en una base incompatible, el resultado pierde correlación con el estado preparado. Esta incompatibilidad introduce incertidumbre y perturbación ante mediciones no autorizadas, por lo que constituye uno de los mecanismos que permite detectar espionaje \citep{scarani2009}.

La seguridad básica se apoya en tres hechos. Primero, como demostraron \cite{wootters1982}, los estados cuánticos desconocidos no pueden copiarse perfectamente, principio conocido como teorema de no clonación. Segundo, medir en una base incompatible introduce perturbaciones estadísticas observables por Alicia y Bob \citep{bennett1984,nielsen2010}. Tercero, la información parcial que pudiera haber obtenido Eva puede reducirse mediante amplificación de privacidad, sacrificando parte de la clave reconciliada para producir una clave final más corta y más segura \citep{gisin2002,scarani2009}.

En una ejecución típica participan Alicia, que prepara estados, y Bob, que mide. Tras muchas rondas, ambos anuncian por el canal clásico parte de la información no secreta, por ejemplo las bases usadas \citep{bennett1984,gisin2002}. A partir de ahí calculan la tasa de error cuántico, o QBER, que mide la fracción de bits discrepantes entre Alicia y Bob en una muestra de la clave \citep{gisin2002,scarani2009}. Si el error supera un umbral, abortan porque el canal puede estar demasiado degradado o intervenido; si es aceptable, aplican reconciliación de información y amplificación de privacidad \citep{scarani2009,lo2014}.

La reconciliación de información es el procedimiento mediante el cual Alicia y Bob corrigen diferencias entre sus cadenas de bits sin revelar completamente la clave \citep{gisin2002,scarani2009}. Después, como señalan \cite{scarani2009}, la amplificación de privacidad transforma la clave corregida en una clave más corta sobre la que Eva tiene información despreciable dentro del modelo de seguridad considerado. Así, la clave final no procede de una única medición perfecta, sino de un procedimiento estadístico sobre muchas rondas.

También es importante distinguir entre seguridad teórica y seguridad de implementación. Un protocolo puede ser seguro en un modelo ideal, pero un sistema real incluye pérdidas, ruido, fuentes imperfectas, detectores no ideales y posibles canales laterales. Por ello, los protocolos modernos no solo especifican cómo preparar y medir cúbits, sino también qué supuestos se hacen sobre los dispositivos y qué contramedidas se necesitan frente a ataques prácticos \citep{scarani2009,lo2014}.